% Generated by maint/texbuild/theory_tex.py from theory/thought_experiments/engine/hash_boundary_functional_folding.md.
% Edits here are overwritten. Edit the markdown, then run the script again.
\chapter{\texorpdfstring{The boundary functional of a hash keyspace, and folding}{The boundary functional of a hash keyspace, and folding}}
\label{chap:hash-boundary-functional-folding}

\section{\texorpdfstring{The question}{The question}}

What would the boundary functional be for a given hash, with at most a 30 digit 0-9a-fA-F keyspace?

\section{\texorpdfstring{1. The scale of a 30 digit hex keyspace}{1. The scale of a 30 digit hex keyspace}}

A 30 digit string using characters 0-9 and a-f (or A-F) operates in base 16. Each hex digit represents 4 bits of information.

\begin{itemize}
\item \textbf{Cardinality:} a 30 digit hex string yields a keyspace of \(16^{30} = (2^4)^{30} = 2^{120}\) possible values.
\item \textbf{Comparison:} smaller than SHA-256's native output (\(2^{256}\)), a \(2^{120}\) keyspace is still vast, roughly \(1.32 \times 10^{36}\) combinations. Even at trillions of evaluations per second, the space remains physically bounded.
\end{itemize}

\section{\texorpdfstring{2. The boundary functional Φ}{2. The boundary functional Φ}}

The boundary functional is the characteristic function that defines the exact edge between a failed search and a hit (the “floor”). It maps a hash output \(x\) in the \(2^{120}\) space to a binary decision or a distance. With \(T\) the target threshold:

\[
\Phi(x) = \begin{cases} 1 & \text{if } x < T \\ 0 & \text{if } x \ge T \end{cases}
\]

Or, as a measure of how close a search came to the boundary:

\[
\Phi(x) = |x - T|
\]

\section{\texorpdfstring{3. The boundary functional in a single cycle sweep}{3. The boundary functional in a single cycle sweep}}

\begin{itemize}
\item \textbf{Simultaneous evaluation:} the grid computes the hash for every nonce in the array concurrently, producing a vector of outputs \([x_1, x_2, \ldots, x_N]\).
\item \textbf{Parallel functional application:} \(\Phi\) is applied to every element at once.
\item \textbf{The latch (the floor):} a reduction tree, or a recursive fold, scans the resulting binary vector; the first 1 trips the latch and locks that index as the floor of the search.
\end{itemize}

Within a 30 digit hex keyspace (\(2^{120}\) states) the boundary functional is the threshold condition, \(\Phi(x) = 1\) when \(x < T\), whether it is evaluated sequentially through a tower recursion or across a spatial array at once.

\section{\texorpdfstring{4. Folding: even with a leading 80 required, there is no dimensional restriction}{4. Folding: even with a leading 80 required, there is no dimensional restriction}}

Folding schemes (Nova, SuperNova, HyperNova) remove the growth that recursion otherwise causes in verifiable computation.

\begin{itemize}
\item \textbf{Deferred verification:} instead of proving the inner logic recursively at every step, a folding scheme defers the expensive final check.
\item \textbf{Instance compression:} two constraint satisfaction instances (step \(i\) and step \(i+1\)) fold into one instance of the same size.
\item \textbf{Constant circuit size:} whether 10 nonces or \(10^{10}\) are folded, the verification circuit stays the same size.
\end{itemize}

With a severe boundary condition, such as a long required prefix:

\begin{itemize}
\item \textbf{The stream:} the nonces are processed as a continuous stream.
\item \textbf{The accumulator:} the folding scheme compresses the running state into a compact vector.
\item \textbf{The final argument:} one succinct argument at the end proves that the accumulated trace satisfied the boundary functional.
\end{itemize}
